\section{Introduction}
\label{sec:introduction}

Large Language Models (LLMs) are undergoing a fundamental transformation~\citep{bommasani2021opportunities,min2023recent,naveed2025comprehensive,wang2024survey,xi2023rise,qin2025largelanguagemodelsmeet}. What began as statistical language generation has expanded into AI systems that can reason, act, remember, and complete tasks in open-ended digital environments~\citep{yao2022react,park2023generative,zhang2025survey}. Early progress was driven by scaling autoregressive Transformers and instruction-aligned chat interfaces, enabling systems to compress broad world knowledge into fluent responses~\citep{vaswani2017attention,brown2020language,kaplan2020scaling,hoffmann2022chinchilla,achiam2023gpt}. More recently, the frontier has shifted toward models that deliberate over difficult problems, invoke tools, interact with environments, and coordinate multi-step workflows~\citep{wei2022chain,deepseek2025r1,schick2023toolformer,qin2023toolllm,wu2023autogen,hong2023metagpt}. The central question is therefore \textit{\textbf{no longer limited to}} \textit{how can a model generate a better answer?} \textit{\textbf{Instead}}, it is how \textit{how can an AI system reliably transform user intent into completed work?} This redefines the human-AI relationship, marking the shift \textbf{from Chatbot to Digital Colleague}~\citep{openai2022chatgpt,bigscience2022bloom,google2024gemma2,jiang2023mistral7b,chen2023palix}.

This survey organizes and analyzes persistent autonomous AI along two tightly coupled dimensions. 
\ding{172} \textit{\textbf{The first dimension concerns the cognitive core}}. It explores how models generate, understand, and reason, spanning two eras, shown in Figure~\ref{fig:map}'s Chatbot and Thinking LLM panels. In the \textit{Chatbot Era}, LLMs behave like fast ``System-1'' generators: they compress parametric knowledge and produce fluent responses, but struggle with deep reasoning, verification, and long-horizon consistency~\citep{petroni2019language,huang2023towards,dziri2023faith,valmeekam2022large}. In the \textit{Thinking LLM Era}, models increasingly leverage inference-time computation, Chain-of-Thought prompting, reflection, process supervision, and reinforcement learning to support slower, more deliberate, and more reliable problem-solving processes~\citep{snell2024scaling,kojima2022large,madaan2023selfrefine,shinn2023reflexion,lightman2024lets,shao2024deepseekmath,deepseek2025r1}. PartI traces the transition to slow, reasoning-centered cognition\citep{schneider2025generative,wei2026agentic,patil2025advancing,zhao2026edge,hu2024survey}.


\ding{173} \textit{\textbf{The second dimension concerns tool-augmented task execution}}. This dimension asks how a stronger cognitive core becomes a system that can act in dynamic and complex external environments, and it also contains two eras, as illustrated in the Agent and OpenClaw panels of Figure~\ref{fig:map}. In the \textit{Agent Era}, LLMs move from passive dialogue systems into active systems that call APIs, browse websites, write code, manipulate files, and collaborate with other agents~\citep{schick2023toolformer,qin2023toolllm,yao2022react,wu2023autogen,hong2023metagpt,xi2023rise,wang2024survey}. Yet, these early agents still remain highly fragile: incorrect action formats, missing observations, failed tool calls, or unrecovered intermediate errors can derail the entire trajectory. In the \textit{OpenClaw Era}, tool use is embedded into persistent workspaces with files, terminals, browsers, logs, permissions, reusable skills, and verification procedures, enabling agents to maintain context, monitor progress, recover from failures, and verify final workspace states~\citep{xu2025toward,guo2024large,lei2025large,sun2025survey,plaat2025agentic}.

Within this two-dimensional framework, the key thesis of this survey is that \textbf{Workspace + Skill} provides the mechanism that turns chatbot-style interaction into durable digital-colleague work~\citep{wang2024openhands,yang2024sweagent,drouin2024workarena}. A \textbf{Workspace} is a persistent digital environment for AI operations, including files, terminals, browsers, editors, repositories, calendars, documents, databases, and domain-specific applications~\citep{xie2024osworld,zhou2023webarena,drouin2024workarena,wang2024openhands}. A \textbf{Skill} is a reusable, parameterizable procedure for completing tasks, including planning, tool sequencing, intermediate checks, error recovery, and validation~\citep{wang2023voyager,shinn2023reflexion,madaan2023selfrefine}. Together, they move LLMs beyond episodic responses and atomic tool calls: the workspace provides state, memory, evidence, and consequences, while the skill provides reusable operational knowledge~\citep{schick2023toolformer,qin2023toolllm,wang2023voyager,wang2024openhands}.

This perspective reframes data and evaluation paradigms in LLM development. For chatbots, data is often organized as instruction-response pairs and evaluation measures final-answer correctness or human preference. For reasoning models, data includes long-form Chain-of-Thought traces, process supervision, and verifiable rewards, while evaluation expands toward reasoning-process judgment. For agents and workspace-based systems, the fundamental unit of learning becomes the \textbf{state--action--observation trajectory}, and evaluation shifts from answer quality to task closure—whether the system reaches the intended final state under reproducible, auditable, and safe conditions~\citep{luo2025large,maestre2024beyond,du2026survey,barua2024exploring,yao2025survey}.

Despite impressive progress, current systems face major structural bottlenecks~\citep{liu2023agentbench,zhou2023webarena,jimenez2023swebench,xie2024osworld}. Reasoning can remain ungrounded or hallucinated during factual verification~\citep{huang2023hallucination,ji2023hallucination}; long-horizon execution is brittle as errors accumulate across toolchains\citep{liu2023agentbench,zhou2023webarena,jimenez2023swebench}; memory and state management often depend on transient context windows\citep{packer2023memgpt,zhang2025survey}; and safety becomes harder when outputs are executable actions with side effects~\citep{ruan2024identifying,debenedetti2024agentdojo,zhan2024injecagent}. These challenges highlight that the transition from chatbot to digital colleague requires both stronger foundation models and better execution substrates, skill abstractions, evaluation environments, and governance mechanisms~\citep{wang2024openhands,yang2024sweagent,wang2023voyager,drouin2024workarena,li2026prism,zhao2026clawguard}.

Accordingly, this survey reviews the field through four parts. Part~I examines the evolution of the LLM cognitive core, from chatbot-era language generation to thinking LLMs driven by long reasoning chains and reinforcement-learning-based cognition. Part~II studies tool-augmented task execution, from early agents to OpenClaw-style systems oriented toward workspace intelligence, skill-based execution, reliability, and governance. Part~III explains why \textbf{Workspace + Skill} is a decisive leap from ephemeral interactions to persistent stateful work and from ad-hoc prompts to composable capability packages. Part~IV analyzes the accompanying shifts in data and evaluation, from knowledge corpora and instruction data to action trajectories, process verification, and task-closure-oriented benchmarks. We then discuss open challenges and future directions toward reliable, self-evolving AI ecosystems.

The main contributions of this survey are summarized as follows:
\begin{itemize}[leftmargin=*, itemsep=2pt, topsep=3pt]
    \item \textbf{A two-dimensional view of AI:} We organize this evolution along two complementary dimensions: cognitive-core evolution (Chatbot and Thinking LLM) and tool-augmented task execution (Agent and OpenClaw), with Workspace + Skill as the mechanism for the Chatbot to Digital Colleague.
    \item \textbf{A unified account of Workspace + Skill, data, and evaluation:} We identify workspace persistence and reusable skills as mechanisms for task completion, and connect this shift with the move from instruction-response pairs to state--action--observation trajectories and task-closure evaluation.
    \item \textbf{A socio-technical roadmap for reliable autonomous AI systems:} We summarize challenges in long-horizon reliability, memory, safety, governance, skills, and system maintenance. We also discuss the broader implications of Digital Colleague systems for human–AI collaboration, including questions of ethics, skills, work pace, creativity, privacy, and asset boundaries.
\end{itemize}

Overall, this survey clarifies how LLMs are moving from conversational chatbots toward dependable digital colleagues. Importantly, the next generation of generative AI will be defined by self-evolving systems: integrated ecosystems in which models, workspaces, tools, skills, memories, evaluators, and governance mechanisms continuously convert operational experience into reusable skills, updated memories, stronger verification signals, safer policies, and more reliable work outcomes.
